\chapter{Orientation spectral no-go: точная изоспектральность двух ветвей}

Common up--down loop оставил два minima
\begin{equation}
 \mathcal B_-:
 \left(W=-\frac12,\theta=0\right),
 \qquad
 \mathcal B_+:
 \left(W=+\frac12,\theta=\pi\right).
\end{equation}
Проверим, может ли иной spectral moment или иная вещественная cutoff function
различить их без добавления нового field content.

\section{Characteristic-polynomial test}

Для первичного \(12\times12\) family operator обе ветви имеют один и тот же
characteristic polynomial:
\begin{multline}
 \chi_-(x)=\chi_+(x)\\
 =
 \frac14
 (x^2-x-1)(x^2+x-1)
 (2x^4-2x^3-6x^2+4x+3)\\
 {}\times
 (2x^4+2x^3-6x^2-4x+3).
\end{multline}
Characteristic polynomials \(D_-^2\) и \(D_+^2\) также совпадают. Поэтому
\begin{equation}
 \Tr f(D_-^2)=\Tr f(D_+^2)
\end{equation}
для любой функции \(f\), определённой на их конечном spectrum.

\begin{theorem}[All-orders real spectral no-go]
Две оставшиеся noncommuting orientation branches точно изоспектральны.
Следовательно, никакая замена cutoff profile, добавление высших even moments
или полный exact functional \(\Tr f(D_F^2)\) не снимают их вырождение.
\end{theorem}

Машинно проверены все power traces до \(D^{12}\), что вместе с равенством
characteristic polynomials фиксирует все последующие moments по
Cayley--Hamilton.

\section{Grading и determinant также не помогают}

Для bipartite grading \(\gamma\) прямой расчёт даёт
\begin{equation}
 \Tr(\gamma D_\pm^{2n})=0,
 \qquad n=1,\ldots,6.
\end{equation}
Кроме того,
\begin{equation}
 \det D_-=\det D_+=\frac94.
\end{equation}
Ориентационно-нечётный loop observable также исчезает в обоих minima:
\begin{equation}
 W\sin\theta=0.
\end{equation}

Таким образом, ни ordinary trace, ни supertrace, ни determinant magnitude не
создают selector в текущем finite menu.

\section{Минимальная развилка после no-go}

Ориентация может появиться только из структуры, чувствительной не к spectrum
\(D^2\), а к фазе fermionic measure или к ориентированному boundary datum.
Остаются три принципиальных маршрута:
\begin{enumerate}
\item Pfaffian phase или \(\eta\)-invariant chiral fermion operator;
\item квантованный boundary Chern--Simons/anomaly-inflow term;
\item явный выход из класса вещественного even spectral action.
\end{enumerate}

\begin{nogocriterion}[Запрет spectral-profile rescue]
После доказанной изоспектральности запрещено искать orientation selector
подбором новой функции \(f\), коэффициентов \(\Tr D^{2n}\) или иной
нормировки bosonic trace. Такой поиск математически не может разделить
\(\mathcal B_-\) и \(\mathcal B_+\).
\end{nogocriterion}

Следующий гейт должен начинаться с fermionic Pfaffian/\(\eta\)-phase и
проверить, не сокращается ли она между \(J\)-сопряжёнными rectangles.
Exact characteristic и moment ledger сохранён в
\path{s2t_v4_orientation_spectral_no_go_results.json}.